% Part: computability
% Chapter: recursive-functions
% Section: introduction

\documentclass[../../../include/open-logic-section]{subfiles}

\begin{document}

\olfileid{cmp}{rec}{int}
\olsection{પરિચય}

સંગણનીયતાનો ગણિતીય સિદ્ધાંત વિકસાવવા માટે સૌપ્રથમ સંગણનીયતાનું એક
\emph{નિદર્શ} ઘડવું પડે. આજે આપણે સંગણનીયતાને કમ્પ્યુટર જે કામ કરે છે
તેના પ્રકાર તરીકે સમજીએ છીએ, અને કમ્પ્યુટર પ્રતીકો સાથે કામ કરે છે. પરંતુ
સંગણનીયતાના સિદ્ધાંતોના વિકાસની શરૂઆતમાં સંગણનાનું આદર્શ ઉદાહરણ
\emph{સંખ્યાત્મક} ગણતરી હતું. ગણિતશાસ્ત્રીઓને હંમેશાં સંખ્યા-સિદ્ધાંતનાં
સંગણન કરી શકાય એવાં વિધેયોમાં, એટલે કે $f\colon \Nat^n \to \Nat$
પ્રકારનાં વિધેયોમાં, રસ રહ્યો છે. તેથી સંગણનીયતાના સિદ્ધાંતની શરૂઆતમાં
આવાં વિધેયોનો અભ્યાસ થયો તે આશ્ચર્યજનક નથી. પ્રાકૃતિક સંખ્યાઓના સરવાળા,
ગુણાકાર અને ઘાતાંક જેવા સૌથી પરિચિત સંગણનીય સંખ્યાત્મક વિધેયોમાં એક
રસપ્રદ સામ્ય છે: તેમને \emph{પુનરાવર્તી રીતે} વ્યાખ્યાયિત કરી શકાય છે.
આથી પુનરાવર્તી વ્યાખ્યાઓને આધારે \emph{સંગણનીય વિધેય}ની સામાન્ય વ્યાખ્યા
આપવાનો પ્રયાસ સ્વાભાવિક છે. સંખ્યા-સિદ્ધાંતનાં વિધેયોને પુનરાવર્તી રીતે
વ્યાખ્યાયિત કરવાની ઘણી રીતોમાંથી અહીં એક ખાસ સરળ પદ્ધતિ કેન્દ્રસ્થાને
છે: કહેવાતું \emph{આદિમ પુનરાવર્તન}.

સંગણનીય વિધેયો ઉપરાંત સંગણનીય ગણો અને સંબંધોમાં પણ આપણને રસ હોઈ શકે.
આપેલી સંખ્યા ગણનો !!a{element} છે કે નહીં, તેનો જવાબ સંગણનથી મેળવી
શકાય તો તે ગણ સંગણનીય છે. તેવી જ રીતે, $\tuple{n_1, \dots, n_k}$
બહુજોડ સંબંધનો !!a{element} છે કે નહીં, તે સંગણનથી નક્કી કરી શકાય તો
તે સંબંધ સંગણનીય છે. ગણ અથવા સંબંધના \emph{લક્ષણ વિધેય}ને ધ્યાનમાં
લેવાથી, સંગણનીય ગણો અને સંબંધોની ચર્ચા સંગણનીય વિધેયોની ચર્ચામાં
સમાવી શકાય છે. તેથી આદિમ પુનરાવર્તી સંબંધો પણ વ્યાખ્યાયિત કરી શકાય;
દા.ત., ``$n$ વડે $m$ શેષ વિના ભાગાય છે'' એવો સંબંધ આદિમ પુનરાવર્તી છે.

તેમ છતાં, આદિમ પુનરાવર્તી વિધેયો---માત્ર આદિમ પુનરાવર્તનથી વ્યાખ્યાયિત
થતાં વિધેયો---જ બધા સંગણનીય સંખ્યા-સિદ્ધાંતના વિધેયો નથી. આદિમ
પુનરાવર્તનના અનેક સામાન્યકરણો વિચારવામાં આવ્યા છે. પરંતુ વિધેયોનું
સંગણન કરવા માટે સૌથી શક્તિશાળી અને વ્યાપકપણે સ્વીકારાયેલી વધારાની રીત
અનિયત મર્યાદા સુધીની શોધ છે. તેનાથી \emph{આંશિક પુનરાવર્તી વિધેયો}ની
વ્યાખ્યા મળે છે અને તેની સાથે સંકળાયેલી \emph{સામાન્ય પુનરાવર્તી
વિધેયો}ની વ્યાખ્યા પણ મળે છે. સામાન્ય પુનરાવર્તી વિધેયો સંગણનીય છે
અને દરેક ઇનપુટ માટે વ્યાખ્યાયિત છે; તેમની વ્યાખ્યા એવા આંશિક
પુનરાવર્તી વિધેયોને ચોક્કસ રીતે ઓળખાવે છે જે વાસ્તવમાં દરેક ઇનપુટ માટે
વ્યાખ્યાયિત હોય. પુનરાવર્તી વિધેયો સંગણનનાં બીજા દરેક નિદર્શનું
(ટ્યુરિંગ યંત્રો, લૅમ્બડા કલન વગેરેનું) અનુકરણ કરી શકે છે. તેથી તેઓ
સંગણનનાં સ્વીકૃત અનેક નિદર્શોમાંનું એક નિદર્શ છે.

\end{document}
